\paragraph*{Problems}

\begin{enumerate}[label=\thesection.\arabic*.,ref=\thesection.\arabic*,start=1,series=book-problems]

\item\label{ex:bump-is-test-function} 证明 $\phi_{\bm{y},\epsilon}\in\mathcal{D}(\mathbb{R}^m)$.

\item\label{ex:test-function-support-escape} 证明序列 $\phi_n(\bm{x})=n^{-1}\phi_{0,\epsilon}(\bm{x})$ 在 $\mathcal{D}(\mathbb{R}^m)$ 中收敛到零. 证明序列 $\psi_n(\bm{x})=n^{-1}\phi_{0,\epsilon}(\bm{x}/n)$ 一致收敛到零, 且所有导数也如此. 为什么 $\psi_n$ 却不在 $\mathcal{D}(\mathbb{R}^m)$ 中收敛到零? 它在 $\mathcal{S}(\mathbb{R}^m)$ 中收敛到零吗?

\item\label{ex:prove-test-function-uniform-approximation} 证明 Theorem~\ref{thm:test-functions-uniform-approximation}.

\item\label{ex:regular-distributions-injective} 设 $f$ 和 $g$ 为 $C(\Omega)$ 中两个不同的函数. 证明它们作为广义函数也不相同.

\item\label{ex:delta-not-continuous-function} 证明 Dirac delta 函数不能与任何连续函数等同.

\item\label{ex:distribution-periodicity-radial-symmetry} 解释广义函数的周期性或径向对称性意味着什么.

\item\label{ex:distribution-local-vanishing} 设 $f$ 为 $\Omega$ 上的广义函数, 且 $G$ 为 $\Omega$ 的开子集. 假设 $G$ 中每一点都有一个使 $f$ 为零的邻域. 证明 $f$ 在 $G$ 上为零. Hint: 使用单位分解论证.

\item\label{ex:pairing-away-from-distribution-support} 证明, 若 $\phi$ 在 $f$ 的支集的某邻域内为零, 则 $(f,\phi)=0$. 若 $\phi$ 只在 $f$ 的支集上为零, 是否足够?

\item\label{ex:principal-value-regularization} 证明 \eqref{eq:principal-value-one-over-x-regularization} 确实定义了一个广义函数, 且这一定义不依赖于 $\epsilon$. 如何定义 $1/x^2$ 的一个正则化?

\item\label{ex:compactly-supported-distribution-density} 证明每个分布都是一个紧支集分布序列的极限. Hint: 令 $f_n=f\psi_n$, 其中 $\psi_n$ 为 $C^\infty$ 截断函数.

\item\label{ex:distribution-multiplication-not-continuous} 证明在 $\mathcal{D}^{\prime}(\mathbb{R})$ 中 \[\lim_{n\to\infty}\sin(nx)=0\] 但 \[\lim_{n\to\infty}\sin^2(nx)\ne0.\] 因而, 即使在分布乘法有定义的地方, 它也不是连续运算.

\item\label{ex:delta-approximation-square-divergence} 设 $f_n$ 为 \eqref{eq:delta-approximating-rectangular-sequence} 定义的序列. 证明 \[\lim_{n\to\infty}f_n^2\] 在分布意义下不存在. 然而, 证明 \[\lim_{n\to\infty}f_n^2-n\delta\] 存在.

\item\label{ex:gaussian-distribution-limit} 求分布意义下的极限 \[\lim_{n\to\infty}\sqrt n\exp(-nx^2)\]

\item\label{ex:distribution-not-tempered-convergence} 给出 $\mathcal{S}^{\prime}(\mathbb{R})$ 中的一个序列, 它在 $\mathcal{D}^{\prime}(\mathbb{R})$ 中收敛到零, 但不在 $\mathcal{S}^{\prime}(\mathbb{R})$ 中收敛到零.

\item\label{ex:rapid-tests-weighted-uniform-convergence} 证明序列 $\phi_n$ 在 $\mathcal{S}(\mathbb{R}^m)$ 中收敛, 当且仅当对每个 $k\in\mathbb{N}\cup\{0\}$ 和每个 $\alpha$, $|\bm{x}|^kD^\alpha\phi_n(\bm{x})$ 均一致收敛.

\item\label{ex:tempered-distributions-sequential-completeness} 证明 $\mathcal{S}^{\prime}(\mathbb{R}^m)$ 是序列完备的.

\item\label{ex:distribution-gluing-open-cover} 设 $U_i$, $i\in\mathbb{N}$ 为开集, 满足 $\Omega=\bigcup_{i\in\mathbb{N}}U_i$. 给定 $f_i\in\mathcal{D}^{\prime}(U_i)$, 使 $f_i$ 与 $f_j$ 在 $U_i\cap U_j$ 上一致. 证明存在 $f\in\mathcal{D}^{\prime}(\Omega)$, 使在 $U_i$ 上 $f=f_i$.

\item\label{ex:distribution-test-sequence-pairing} 证明 Lemma~\ref{lem:distribution-test-sequence-pairing}.

\item\label{ex:locally-finite-cover-refinement} (a) 设 $\Omega$ 为 $\mathbb{R}^m$ 的任意开子集. 证明存在具有 Definition~\ref{def:partition-of-unity} 中性质的一族子集.

(b) 设 $\{U_i\}$ 为 $\Omega$ 的任意可数开覆盖. $\{U_i\}$ 的一个加细是由开集 $V_k$ 构成的覆盖, 其中每个 $V_k$ 是某个 $U_i$ 的子集. 给定 $\Omega$ 的任意开覆盖, 证明存在满足 Definition~\ref{def:partition-of-unity} 中性质的加细.

\end{enumerate}
